\section{Introduction}\label{sec:intro}
LLM-based trading agents such as FinMem~\cite{yu2025finmem}, FinCon~\cite{yu2024fincon} and TradingAgents~\cite{xiao2025tradingagents} retrieve news from a memory store, reason over it together with market data, and emit sequential trading decisions~\cite{li2025investorbench,zhang2024finagent}. Their text inputs are an attack surface. Unlike a jailbreak, financial misinformation need not cross any safety boundary: a flipped direction word, an inflated percentage or a swapped cause produces text that still reads as plausible news, is hard to flag without a reference~\cite{jiang2026glisters}, and can change what the agent does. Poisoned retrieval corpora~\cite{zou2025poisonedrag} and injected third-party content~\cite{greshake2023indirect} are established threats to LLM applications; trading agents add a direct financial channel.

AutoRedTrader~\cite{liu2026autoredtrader} showed how far such attacks go. Its attacker combines behavioral-bias prompts, minor semantic edits and style rewrites, filters candidates by similarity and detectability, and reinforces the strategies that changed the agent's decisions. In our replication it flips \val{MInjectedNoneASR}\% of an undefended agent's daily decisions. Detecting the edited text is hard~\cite{jiang2026glisters,wang2026life}, and the attacker filters its candidates against whatever detector it faces. AutoRedTrader also proposes a preliminary countermeasure, time-series-informed grounding, which checks news against recent price windows; we develop it into a complete defense layer.

\textbf{Ground the news in evidence the attacker cannot edit.} The attacker controls what the agent reads, not what the market does. \rhal{} therefore places a defense layer between the agent's policy and its orders (Fig.~\ref{fig:harness}). An evidence model summarizes the market up to day $t$ from trailing price features, microstructure proxies and the realized outcomes of similar past windows. Temporal grounding pulls the agent's text signal toward this evidence in proportion to their conflict. A risk layer then gates and sizes the final decision. Every component uses only data available on day $t$, and tests that mutate future prices and text verify this.

\textbf{Evaluate the defense fairly.} Defense comparisons are easy to get wrong: arms that draw different attacker random numbers differ by attack noise alone, an attacker that adapts to each defense shows each one different documents, and adding \emph{any} documents to the agent's pool changes some decisions even when the text is harmless. We therefore trade an attacked agent and a clean twin in lockstep with shared attacker randomness, replay a recorded attack verbatim against every defense in addition to running the adaptive attacker, and subtract placebo controls so that a defense is credited only for misinformation-driven flips. All conclusions are paired over \val{NseedsRun} simulated markets with bootstrap intervals and multiplicity-corrected tests.

\textbf{Test it where it could fail.} We compare grounding with detection filtering, robust aggregation and an agent that ignores the news. We then replace the lexicon reader with a local LLM, trade real BTC-USD prices, give the attacker the market evidence the defense uses, and let the misinformation move prices.

\textbf{Contributions.}
\begin{itemize}
\item The \emph{RobustHFT defense layer}: a market-evidence model with leak-free historical analogs, temporal grounding of the agent's text signal, and a risk layer that can only reduce exposure. It wraps lexicon and LLM readers alike (Sec.~\ref{sec:defense}).
\item An evaluation protocol for such defenses: a paired clean twin, common random numbers, adaptive and frozen-replay attacks, defense-aware attackers, a price-impact stress test, and placebo controls that separate retrieval noise, bias prompts and misinformation content (Sec.~\ref{sec:protocol}--\ref{sec:setup}).
\item Results over \val{NseedsRun} markets (Sec.~\ref{sec:results}). Grounding cuts misinformation-driven ASR by about \val{RelCutGroundedNet}\% for \val{CmpCostGrounded}\pp{} of clean-market return and beats the standard defenses. The gain holds with an LLM reader, on real prices, against an evidence-aligned attacker and under price impact. A hold-on-conflict rule backfires and a shrink rule removes the penalty, while a conflict-seeking attacker remains the main open threat.
\end{itemize}
